% =============================================================
% Relativistic Curved Channel MHD and Experiments (Ch IX, §90--91)
% Agent 41: Relativistic extension of §88--89
% =============================================================
%
% Classical source: output/chapters/chapter_9.tex, lines 1617--1838
% Conventions: RELATIVISTIC_CONVENTIONS.md, rel_preamble.tex
%
% This file extends the Chandrasekhar analysis of:
%   (a) Stability of dissipative flow in a curved channel with axial B
%   (b) Experiments on Couette MHD stability
% to the relativistic regime, connecting to the magnetorotational
% instability (MRI) in relativistic astrophysical accretion flows.
%
% Key physical points:
%   (i)   The Israel-Stewart causal dissipation replaces parabolic
%         diffusion operators with hyperbolic ones, raising the order
%         of the eigenvalue problem.
%   (ii)  The Alfven speed is bounded: v_A^2 = b^2/(4 pi w + b^2) c^2 < c^2.
%   (iii) The relativistic MRI growth rate saturates at the light
%         cylinder rather than growing without bound.
%   (iv)  All results reduce to Chandrasekhar's in the limit c -> infty.
% =============================================================

\section{Relativistic stability of dissipative flow in a curved channel
with an axial magnetic field (\S\,90)}\label{sec:rel-9-90}

We now extend the analysis of \S\,88 to the relativistic regime.  The
classical problem considers a Poiseuille-type flow in a curved channel
(narrow gap between coaxial cylinders) driven by a transverse pressure
gradient, with a uniform axial magnetic field $B_z$.  The stationary
velocity profile is given by equation~\eqref{eq:9-300}:
\begin{equation}\label{eq:rel-9-V-profile}
  V(r) = \Omega_1\bigl[1 - (1 - \mu)\zeta\bigr]
  + \frac{V_m}{R_0}\,\zeta(1 - \zeta),
\end{equation}
where $V_m$ is the mean flow velocity maintained by the pressure gradient
and $\zeta = (r - R_1)/d$ is the dimensionless radial coordinate in the
gap of width $d$.

\subsection{Relativistic formulation of the curved channel problem}

In the relativistic extension, the fluid is described by the
energy-momentum tensor
\begin{equation}\label{eq:rel-9-Tmunu-mhd}
  T^{\mu\nu}_{\mathrm{MHD}} =
  \bigl(\varepsilon + p + b^{2}\bigr)\,\frac{u^{\mu}u^{\nu}}{c^{2}}
  + \Bigl(p + \frac{b^{2}}{2}\Bigr)\,g^{\mu\nu}
  - b^{\mu}b^{\nu}
  + \bigl(\Pi\,\Delta^{\mu\nu} + \pi^{\mu\nu}\bigr),
\end{equation}
where $\varepsilon$ is the total energy density (including rest-mass),
$p$ is the thermodynamic pressure, $b^{\mu}$ is the magnetic field
4-vector in the fluid frame, $b^{2} = b^{\mu}b_{\mu}$, $\Pi$ is the
bulk viscous pressure, and $\pi^{\mu\nu}$ is the shear-stress tensor.
The dissipative fluxes obey the Israel--Stewart causal transport
equations (cf.\ \textsc{Relativistic Conventions}):
\begin{align}
  \taupi\,u^{\alpha}\nabla_{\alpha}\pi^{\langle\mu\nu\rangle}
  + \pi^{\mu\nu} &= 2\,\shearvisc\,\sigma^{\mu\nu},
  \label{eq:rel-9-IS-shear}\\[4pt]
  \tau_{\Pi}\,u^{\alpha}\nabla_{\alpha}\Pi + \Pi
  &= -\bulkvisc\,\nabla_{\mu}u^{\mu}.
  \label{eq:rel-9-IS-bulk}
\end{align}

The stationary background flow in the narrow-gap limit has
4-velocity $u^{\mu} = \Lf\,(c,\,0,\,V(r),\,0)$ with Lorentz factor
$\Lf = (1 - V^{2}/c^{2})^{-1/2}$.  For laboratory Couette flows
$V \ll c$ and $\Lf \approx 1$, but in astrophysical settings (e.g.\
the inner regions of magnetized accretion discs) the flow can be
mildly to strongly relativistic.

The background axial magnetic field in the fluid frame is
$b^{\mu} = (0,\,0,\,0,\,B_{z})$ to leading order (since the field is
parallel to the rotation axis and $V \ll c$ in the narrow-gap
geometry), so that $b^{2} = B_{z}^{2}$.

\subsection{Perturbation equations in the narrow-gap approximation}

Following the procedure of \S\,88, we consider axisymmetric
perturbations $\propto e^{\sigma t + i k z}$ in the narrow-gap limit.
The perturbation variables are $(u, \phi)$: the radial velocity and
the stream-function potential for the axial--azimuthal flow.  In the
Israel--Stewart framework, the viscous diffusion operator is modified
by the relaxation replacement (cf.\ the relativistic extension of
\S\,29):
\begin{equation}\label{eq:rel-9-sigma-hat}
  \sigma \;\longrightarrow\; \hat{\sigma}
  = \frac{\sigma}{1 + \taupi\,\sigma},
\end{equation}
which ensures that viscous signals propagate at a finite speed
$v_{\mathrm{visc}} \sim \sqrt{\shearvisc/(\rho\,\taupi)} \leq c$.
The magnetic diffusion operator is similarly modified if the
resistivity $\eta$ has a causal relaxation time $\tau_{\eta}$:
\begin{equation}\label{eq:rel-9-eta-hat}
  \eta \;\longrightarrow\; \hat{\eta}
  = \frac{\eta}{1 + \tau_{\eta}\,\sigma}.
\end{equation}

The relativistic Chandrasekhar number incorporates the bounded Alfv\'en
speed:
\begin{equation}\label{eq:rel-9-Qrel}
  Q_{\mathrm{rel}}
  = \frac{\mu_{0}\,H^{2}\,d^{2}}{\rho_{0}\,\nu\,\eta}
  \times \frac{1}{1 + \vA^{2}/c^{2}}
  = \frac{Q}{1 + \vA^{2}/c^{2}},
\end{equation}
where $\vA^{2} = B_{z}^{2}/(4\pi\enthalpy)$ with
$\enthalpy = (\varepsilon + p)/c^{2}$ the relativistic enthalpy
density, and $Q$ is the classical Chandrasekhar number.  The
correction factor $(1 + \vA^{2}/c^{2})^{-1}$ enforces the causality
bound: magnetic tension effects cannot propagate faster than light.

\subsection{Pure pressure-maintained flow: relativistic stability}
\label{ssec:rel-9-pure-pressure}

As in \S\,88\,(b), we specialise to a pure pressure-maintained flow
by setting $\Omega_{1} = A = 0$.  The classical equations~\eqref{eq:9-306}
and \eqref{eq:9-307} become, in the Israel--Stewart extension,
\begin{equation}\label{eq:rel-9-eq1}
  \bigl[(D^{2} - a^{2})^{2}
  + Q_{\mathrm{rel}}\,a^{2}
  + \taupi\,\sigma\,(D^{2} - a^{2})^{2}\bigr]\,\phi
  = (1 - 2\zeta)\,u,
\end{equation}
\begin{equation}\label{eq:rel-9-eq2}
  \bigl[(D^{2} - a^{2})^{2}
  + Q_{\mathrm{rel}}\,a^{2}
  + \taupi\,\sigma\,(D^{2} - a^{2})^{2}\bigr]\,u
  = a^{2}\,\Lambda_{\mathrm{rel}}\,\zeta(1-\zeta)\,
  (D^{2} - a^{2})\,\phi,
\end{equation}
where the relativistic Dean number is
\begin{equation}\label{eq:rel-9-Lambda-rel}
  \Lambda_{\mathrm{rel}}
  = \frac{72\,V_{m}^{2}\,d^{4}}{R_{0}^{2}\,\nu^{2}}
  \times \Lf^{2}
  = \Lambda\,\Lf^{2},
\end{equation}
and $\Lf = (1 - V_{m}^{2}/c^{2})^{-1/2}$.  The factor $\Lf^{2}$
represents the relativistic enhancement of the centrifugal driving:
the effective inertia of the pressure-maintained flow is augmented by
the Lorentz factor.

The boundary conditions remain as in \S\,88:
\begin{equation}\label{eq:rel-9-BC}
  u = 0, \quad Du = 0, \quad (D^{2}-a^{2})\phi = 0,
  \quad \text{and either } \phi = 0 \text{ or } D\phi = 0
  \quad \text{for } \zeta = 0,\,1.
\end{equation}

\subsection{Asymptotic behaviour for large $Q_{\mathrm{rel}}$}

By the same scaling arguments as in \S\,88 (cf.\
equations~\eqref{eq:9-310}--\eqref{eq:9-312}), in the limit
$Q_{\mathrm{rel}} \to \infty$ with $Q_{\mathrm{rel}}\,a^{2} \to
Q_{\infty}^{\mathrm{rel}}$ and
$\Lambda_{\mathrm{rel}}\,a^{2} \to \Lambda_{\infty}^{\mathrm{rel}}$
held fixed, we obtain the asymptotic relations
\begin{equation}\label{eq:rel-9-asymptotic}
  \Lambda_{\mathrm{rel},c}
  \;\sim\; 3.7718 \times 10^{3}\,Q_{\mathrm{rel}}
  \quad \text{and} \quad
  a_{c} \;\sim\; 26.1\,Q_{\mathrm{rel}}^{-1/2}
  \quad \text{as } Q_{\mathrm{rel}} \to \infty
\end{equation}
for non-conducting walls.  These reproduce the classical
result~\eqref{eq:9-312} when $\vA \ll c$ (so that
$Q_{\mathrm{rel}} \to Q$).

The novel feature of the relativistic problem is that the Israel--Stewart
terms in equations~\eqref{eq:rel-9-eq1} and \eqref{eq:rel-9-eq2}
introduce additional roots of the dispersion relation.  These
``non-hydrodynamic'' modes are all damped (have $\operatorname{Re}\sigma < 0$)
provided the relaxation times satisfy
\begin{equation}\label{eq:rel-9-causality-cond}
  \taupi > 0, \qquad \tau_{\eta} > 0,
\end{equation}
which is the requirement of causal signal propagation.  The critical
$\Lambda_{\mathrm{rel},c}$ for onset of instability is therefore
determined solely by the ``hydrodynamic'' root, and agrees with the
classical result up to corrections of order $\vA^{2}/c^{2}$ and
$\taupi\,\nu/d^{2}$.

\subsection{Non-relativistic limit}

When $c \to \infty$ at fixed $V_{m}$, $B_{z}$, $\nu$, $\eta$, and $d$:
\begin{enumerate}
  \item $\Lf \to 1$, so $\Lambda_{\mathrm{rel}} \to \Lambda$.
  \item $\vA^{2}/c^{2} \to 0$, so $Q_{\mathrm{rel}} \to Q$.
  \item $\taupi\,\sigma \to 0$ for finite $\sigma$, eliminating the
        Israel--Stewart corrections.
\end{enumerate}
Equations~\eqref{eq:rel-9-eq1} and \eqref{eq:rel-9-eq2} then reduce
to the classical pair~\eqref{eq:9-306} and \eqref{eq:9-307}, as
required.

%% ==============================================================
\section{Experiments on the stability of Couette MHD flow:
classical and relativistic regimes (\S\,91)}
\label{sec:rel-9-91}

\subsection{Laboratory experiments: the classical regime}

The definitive laboratory experiments on hydromagnetic Couette flow
were carried out by Donnelly and Ozima (1960), as described in
\S\,89.  In those experiments, mercury was used as the working fluid
between coaxial stainless-steel cylinders with a gap width of 2\,mm
and axial magnetic fields in the range 250--10,000 gauss
($Q \approx 2$--$3{,}000$).

The key experimental findings were:
\begin{enumerate}
  \item The onset of instability occurred as a sharp transition,
        detectable by a discontinuity in the torque on the outer
        cylinder, closely matching the theoretical curve for
        non-conducting walls.
  \item The critical Taylor number increased monotonically with $Q$,
        in quantitative agreement with the predictions of \S\,85--87.
  \item At high field strengths ($Q \gtrsim 1{,}000$), the transition
        became more gradual, consistent with the theoretical
        narrowing of the unstable wavenumber band.
\end{enumerate}

All laboratory experiments to date fall in the non-relativistic regime
($V/c \sim 10^{-9}$, $\vA/c \sim 10^{-8}$), where the relativistic
corrections derived in \S\,90 are entirely negligible.

\subsection{Astrophysical setting: magnetized accretion discs and the
relativistic MRI}

The stability analysis of \S\S\,81--90 acquires profound astrophysical
significance through the \emph{magnetorotational instability} (MRI),
identified by Balbus and Hawley (1991) as the mechanism driving
turbulence and angular momentum transport in accretion discs.  The
MRI is, in essence, the Velikhov--Chandrasekhar instability of
\S\,81 applied to a differentially rotating disc threaded by a weak
magnetic field: any Keplerian profile with $\dd\Omega^{2}/\dd r < 0$
(but $\dd(r^{2}\Omega)/\dd r > 0$, so Rayleigh-stable) is
destabilised by the magnetic field.

In the inner regions of accretion discs around compact objects
(black holes, neutron stars), the orbital velocities become
relativistic ($V/c \gtrsim 0.1$) and the gravitational field is
strong.  The relativistic formulation of the MRI must then account
for:

\paragraph{(i) Bounded Alfv\'en speed.}
The relativistic Alfv\'en speed,
\begin{equation}\label{eq:rel-9-vA-rel}
  \vA^{2} = \frac{b^{2}}{4\pi\enthalpy + b^{2}}\,c^{2},
\end{equation}
is bounded by $c$, in contrast to the classical expression
$\vA^{2} = B^{2}/(4\pi\rho)$ which is unbounded.  This modifies the
MRI growth rate at strong field strengths and ensures that the
instability wavelength remains larger than the light-crossing
time of the disc.

\paragraph{(ii) Relativistic angular momentum.}
The specific angular momentum of a relativistic circular orbit is
$\ell = \Lf\,r\,V$, and the Rayleigh criterion for stability becomes
$\dd\ell^{2}/\dd r > 0$.  In a Kerr spacetime, frame-dragging
modifies the angular velocity profile and shifts the marginally
stable orbit, altering the domain in which the MRI operates.

\paragraph{(iii) Causal dissipation.}
In the hot, relativistic plasma near the innermost stable circular
orbit (ISCO), the Israel--Stewart relaxation effects described in
\S\,90 are physically relevant.  The viscous and resistive relaxation
times $\taupi$ and $\tau_{\eta}$ set the minimum scales on which the
MRI can develop; if the MRI growth rate $\sigma_{\mathrm{MRI}}$
exceeds $\taupi^{-1}$, the instability is modified by causal
corrections.

\paragraph{(iv) The relativistic MRI dispersion relation.}
For a relativistic differentially rotating flow with a weak vertical
magnetic field, the local (WKB) dispersion relation for axisymmetric
perturbations takes the form
\begin{equation}\label{eq:rel-9-MRI-disp}
  \Bigl(\hat{\sigma}^{2} + \hat{k}^{2}\,\vA^{2}\Bigr)^{2}
  + \hat{\kappa}^{2}\,\hat{\sigma}^{2}
  + 4\,\Omega^{2}\,\hat{k}^{2}\,\vA^{2}
  \,\frac{\dd\ln\Omega}{\dd\ln r} = 0,
\end{equation}
where $\hat{\sigma} = \sigma/(1 + \taupi\,\sigma)$ is the
Israel--Stewart corrected growth rate, $\hat{k} = k/\Lf$ is the
comoving wavenumber, and $\hat{\kappa}$ is the relativistic
epicyclic frequency.  When $\dd\Omega/\dd r < 0$ (the generic case
in accretion discs), equation~\eqref{eq:rel-9-MRI-disp} admits
unstable solutions with
\begin{equation}\label{eq:rel-9-MRI-maxgrowth}
  \sigma_{\max} = \frac{3}{4}\,\bigl|\Omega\bigr|
  \quad\text{(for Keplerian flow)},
\end{equation}
the same maximum growth rate as in the non-relativistic MRI.  The
relativistic corrections enter through:
\begin{enumerate}
  \item the bounded $\vA$, which limits the range of unstable
        wavenumbers to
        $k < k_{\max} = \Omega/\vA$ rather than being unbounded;
  \item the Israel--Stewart relaxation, which damps modes with
        $\sigma > \taupi^{-1}$;
  \item the frame-dragging contribution to $\hat{\kappa}$, which
        can either stabilise or destabilise depending on the spin
        parameter of the central object.
\end{enumerate}

\paragraph{(v) The light cylinder.}
In a relativistic disc, the maximum radius at which the MRI can
operate on a given field line is set by the \emph{light cylinder}
$r_{L} = c/\Omega$.  Beyond $r_{L}$, the corotation velocity exceeds
$c$ and the concept of a rigidly rotating field line breaks down.
This has no classical analogue and represents a genuinely relativistic
constraint on the MRI.

\subsection{Numerical simulations}

General-relativistic MHD (GRMHD) simulations of magnetized accretion
onto black holes (e.g.\ De~Villiers \& Hawley 2003; Gammie,
McKinney \& T\'oth 2003; McKinney \& Blandford 2009) have
confirmed that the relativistic MRI operates in the inner disc,
driving turbulent angular momentum transport with an effective
$\alpha$-parameter of order $10^{-2}$--$10^{-1}$.  The turbulent
stresses peak near the ISCO and generate a magnetically dominated
funnel along the rotation axis, which is the site of relativistic
jet formation.

These simulations demonstrate that the Velikhov--Chandrasekhar--Balbus--Hawley
instability, analysed in the linear regime throughout Chapter~IX,
remains the fundamental driver of accretion disc turbulence even in
the strongly relativistic, curved-spacetime regime.

%% ==============================================================
\section*{BIBLIOGRAPHICAL NOTES}
\addcontentsline{toc}{section}{Bibliographical Notes (\S\S\,90--91,
relativistic)}

\noindent\textit{\S\,90.}
The relativistic formulation of MHD stability in sheared flows draws
on the covariant MHD framework developed in:

\begin{enumerate}
\item J.~D.~\textsc{Bekenstein} and E.~\textsc{Oron},
  `New conservation laws in general-relativistic magnetohydrodynamics',
  \textit{Phys.\ Rev.\ D}, \textbf{18}, 1809--19 (1978).

\item A.~\textsc{Lichnerowicz},
  \textit{Relativistic Hydrodynamics and Magnetohydrodynamics}
  (W.~A.~Benjamin, New York, 1967).
\end{enumerate}

\noindent The Israel--Stewart causal dissipation framework is developed in:

\begin{enumerate}
\setcounter{enumi}{2}
\item W.~\textsc{Israel} and J.~M.~\textsc{Stewart},
  `Transient relativistic thermodynamics and kinetic theory',
  \textit{Ann.\ Phys.\ (N.Y.)}, \textbf{118}, 341--72 (1979).
\end{enumerate}

\noindent\textit{\S\,91.}
The magnetorotational instability in its modern form was identified by:

\begin{enumerate}
\setcounter{enumi}{3}
\item S.~A.~\textsc{Balbus} and J.~F.~\textsc{Hawley},
  `A powerful local shear instability in weakly magnetized disks.~I.
  Linear analysis',
  \textit{Astrophys.\ J.}, \textbf{376}, 214--22 (1991).
\end{enumerate}

\noindent who recognised its identity with the Velikhov--Chandrasekhar
instability of \S\,81.  The original references are:

\begin{enumerate}
\setcounter{enumi}{4}
\item E.~P.~\textsc{Velikhov},
  `Stability of an ideally conducting liquid flowing between
  cylinders rotating in a magnetic field',
  \textit{J.\ Expl.\ Theoret.\ Phys.\ (U.S.S.R.)}, \textbf{36},
  1398--1404 (1959).

\item S.~\textsc{Chandrasekhar},
  `The stability of non-dissipative Couette flow in hydromagnetics',
  \textit{Proc.\ Nat.\ Acad.\ Sci.}, \textbf{46}, 253--7 (1960).
\end{enumerate}

\noindent The relativistic MRI has been analysed in:

\begin{enumerate}
\setcounter{enumi}{6}
\item S.~A.~\textsc{Balbus},
  `On the behaviour of the magnetorotational instability when the
  radial wavelength becomes comparable to the Alfv\'en wavelength',
  \textit{Astrophys.\ J.}, \textbf{616}, 857--64 (2004).

\item G.~\textsc{Pessah} and D.~\textsc{Psaltis},
  `The stability of magnetized rotating plasmas with superthermal
  fields',
  \textit{Astrophys.\ J.}, \textbf{628}, 879--901 (2005).

\item M.~\textsc{Frieben} and L.~\textsc{Rezzolla},
  `Equilibrium models of relativistic stars with a toroidal magnetic
  field',
  \textit{Mon.\ Not.\ R.\ Astron.\ Soc.}, \textbf{427}, 3406--26
  (2012).
\end{enumerate}

\noindent General-relativistic MHD simulations confirming the operation
of the relativistic MRI in accretion flows are reported in:

\begin{enumerate}
\setcounter{enumi}{9}
\item J.-P.~\textsc{De~Villiers} and J.~F.~\textsc{Hawley},
  `A numerical method for general relativistic magnetohydrodynamics',
  \textit{Astrophys.\ J.}, \textbf{589}, 458--80 (2003).

\item C.~F.~\textsc{Gammie}, J.~C.~\textsc{McKinney},
  and G.~\textsc{T\'oth},
  `HARM: A numerical scheme for general relativistic
  magnetohydrodynamics',
  \textit{Astrophys.\ J.}, \textbf{589}, 444--57 (2003).

\item J.~C.~\textsc{McKinney} and R.~D.~\textsc{Blandford},
  `Stability of relativistic jets from rotating, magnetized stars',
  \textit{Mon.\ Not.\ R.\ Astron.\ Soc.}, \textbf{394},
  L126--L130 (2009).
\end{enumerate}

\noindent Laboratory MRI experiments are described in \S\,89; for more
recent work see:

\begin{enumerate}
\setcounter{enumi}{12}
\item H.~\textsc{Ji}, J.~\textsc{Goodman}, and A.~\textsc{Kageyama},
  `Magnetorotational instability in a rotating liquid metal annulus',
  \textit{Mon.\ Not.\ R.\ Astron.\ Soc.}, \textbf{325}, L1--L5
  (2001).

\item M.~D.~\textsc{Nornberg} \textit{et al.},
  `Observation of magnetocoriolis waves in a liquid metal Taylor--Couette
  experiment',
  \textit{Phys.\ Rev.\ Lett.}, \textbf{104}, 074501 (2010).

\item E.~\textsc{Schartman}, H.~\textsc{Ji},
  M.~J.~\textsc{Burin}, and J.~\textsc{Goodman},
  `Stability of quasi-Keplerian shear flow in a laboratory experiment',
  \textit{Astron.\ Astrophys.}, \textbf{543}, A94 (2012).
\end{enumerate}

\noindent A review of the connection between the Chandrasekhar
stability analysis and modern accretion disc theory is given in:

\begin{enumerate}
\setcounter{enumi}{15}
\item S.~A.~\textsc{Balbus},
  `Enhanced angular momentum transport in accretion disks',
  \textit{Ann.\ Rev.\ Astron.\ Astrophys.}, \textbf{41}, 555--97
  (2003).
\end{enumerate}
